\chapter[An Ocean without a Master]{The Indian Ocean Had No Master, but Was Never Empty}
\section{A Chinese Shard on an African Shore}
Pick up a palm-sized pottery fragment: a white body with blue designs, edges rounded by sand and seawater like a sucked sweet.

It was excavated at Gedi, a ruined Kenyan coastal city near Mombasa, swallowed by forest. Coral-stone palaces, mosques, and houses flourished between the thirteenth and seventeenth centuries. Among many imported finds, Chinese shards stand out. Made over eight thousand kilometres away, they were packed, shipped, unloaded, and resold before breaking on an East African household floor.

Pause over how strange that is.

Eight thousand kilometres without aeroplanes, container ships, or a postal system. More importantly, no emperor or navy controlled the entire route. China's ruler could not govern the mid-Indian Ocean; Arabian rulers could not govern India; Indian kings could not govern East Africa. The ocean had no single master.

Yet the ceramic arrived, along with thousands more. Almost every ancient coastal site from Somalia to Mozambique yields Chinese shards; wealthy households even embedded whole bowls and dishes in tomb pillars and mosque walls. Ceramics were relayed rather than simply transported, hand to hand and port to port, like a decades-long relay among participants mostly unknown to one another and sometimes without a shared language.

For over a millennium, this masterless ocean was among Earth's busiest, most orderly, and most tolerant waters, never empty. Where did its order come from without an emperor? That is our question.

\section[Malacca: Waiting for the Wind]{Morning in Malacca: Everyone Waits for the Wind}
Enter Malacca around 1500, one southwest-monsoon morning. Today's small Malaysian west-coast city was then Southeast Asia's busiest port.

Before full daylight, ships fill the anchorage. Their origins show in their construction: low, slender, triangular-sailed dhows from Arabia and western India; square-bodied Chinese vessels with sails like bamboo-ribbed fans; shallow-draught boats with raised ends from Java and Sumatra. Shore workers can guess cargo from sail silhouettes.

Smells announce the place before languages do. Pepper dominates riverside warehouses; grains escape worn sacks into mud and crunch underfoot. Sandalwood sweetness, salted fish, tar, hemp rope, and thousands of cooking fires mingle.

The harbourmasters' office is open. Malacca has four shahbandars, as Malay calls them: one for Gujarati traders, one for Bengal and Burma, one for local island commerce, and one for ships from China, Ryukyu, and Vietnam. Different languages, measures, gods, and dispute customs require this division. A harbourmaster must be half translator, half judge, and half tax collector.

A Gujarati merchant inspects cotton cloth from Cambay on India's west coast: thousands of bolts, white, dyed, and printed. Island customers wear it; further traders exchange it for spices. He plans a return cargo of cloves and nutmeg, grown only on a few distant volcanic islands. India multiplies their price several times; Europe sends it sky-high.

Behind him, a Chinese merchant bargains with a local broker over porcelain, silk, and iron pots. An Arab sailor eats dates while squatting on a mooring rope; his ship from Aden carried Arabian horses, the ocean's costliest and most demanding cargo. Many die from the voyage's battering, while survivors fetch fortunes. Swahili men repair seams, having brought East African gold dust and ivory to exchange for cloth and bowls.

Everyone does the same thing: watches the sky.

They watch clouds gathering southwest because one invisible dispatcher commands every ship: the monsoon. Southwest winds for half the year, northeast for the other, regular as a pendulum. Before reversal, load cargo, settle accounts, and say goodbye to temporary families. Many merchants never return home; many marry and settle here, as we will see. Once the wind turns, the anchorage empties within weeks like an ebbing tide.

No king orders arrivals or departures. The wind does.

\section{Who Governs This Sea?}
Set the morning aside and consider a fundamental conflict.

Our familiar history grows on land. Historical maps show coloured blocks: Mughal, Ottoman, a great Ming expanse, tiny Malacca. Each has borders, a capital, rulers, tax officials, and frontier troops. Order seems simple: whoever owns this land supplies its rules.

Malacca's sailors see another map of the same Earth, made of points and lines. Ports are points: Aden, Hormuz, Cambay, Calicut, Malacca, Quanzhou, Kilwa. Routes carry month labels and seasonal directions. Here a Mughal emperor differs little from a village headman: both are land people unable to govern water.

The real question emerges:

\textbf{Why did no master not mean chaos?}

Land logic predicts ungoverned waters as pirates' paradise and merchants' graveyard, where seizure makes ownership and strong ships collect protection money. Order supposedly requires a central monopoly of violence. We learn this repeatedly, in history and the everyday claim that unsupervised study periods become chaos.

Yet the Indian Ocean answered differently. For over a millennium before cannon-bearing Europeans arrived, long-distance trade was broadly peaceful. Merchant vessels generally lacked heavy weapons; pirates existed but did not dominate. Ports competed through fairness, lower taxes, and dry warehouses, not more guns. Spices, ceramics, cotton, horses, and gold circulated for a thousand years, sustaining dozens of ports and millions of livelihoods from East Africa to China.

A deeper question follows: if this network truly flourished, why does our history barely include it? We can name emperors but not a fifteenth-century merchant, list dynastic battles but not describe Malacca's morning. Who decides what deserves memory?

By the chapter's end, you will have your own answer.

\section{Sultans, Merchants, and the Unarmed}
Let the opposing parties speak.

Malacca's sultan has a confident case. Why do ships come here rather than trade on a deserted island? Sheltered anchorage, sound warehouses, dependable harbourmasters, honest scales, and markets where thieves are caught. Security, trust, and measures were built by him and his ancestors, not dropped from heaven. Taxes of a few per cent are justified: not charity, but traders still profit. He must obey rules too. Arbitrary increases or confiscation travel with the next monsoon, sending next year's ships to Aceh or Patani. Port competition reins in the sultan.

A Gujarati cloth merchant offers another logic. He owes no king allegiance. Cambay is home; Aden, Malacca, and Calicut hold his warehouses. Partners trusted with gold-worth cargo for a year's separation live in seven or eight ports. They usually share religion or denomination and meet at the same temple or mosque on festivals. This is risk management, not superstition. Where telephones and courts cannot reach, faith and hometown ties guarantee contracts. Defraud one partner and the whole network refuses future credit shipments. Reputation restrains people better than any navy.

A Swahili sailor cares about wind, stars, reefs, and water. Possibly illiterate, he carries oral sailing instructions: which rising constellation signals departure, which cloud means shelter, which green water warns of reefs. Generations of East African, Arabian, and Indian sailors accumulated this knowledge. Fifteenth-century Arab navigator Ahmad ibn Majid wrote dozens of rhymed navigational poems, with turns and anchorages in memorable verse. Maritime technology meant recitation and memory, not instruments and maps.

Then came newcomers. In 1498, three Portuguese ships under Vasco da Gama rounded the Cape of Good Hope and reached Calicut on India's southwest coast. Their own accounts describe hiring a local pilot for the East Africa--India crossing. Remember that the intruders initially found the doorway through local knowledge.

They brought an alien idea: the ocean must have an owner. Merchantmen needed purchased passes or risked supposedly legal plunder. Uncooperative ports faced bombardment. In 1511, Portugal captured Malacca.

State their strongest case seriously rather than mocking a caricature:

Masterless seas mean nobody responsible for order. Who funds anti-piracy action, shipwreck rescue, or anchorage maintenance? Everyone wants public benefits; nobody wants to pay. You know the problem from classroom cleaning. States solve it ashore by monopolising force, levying taxes, and supplying order. Portugal exported that logic to sea. Europeans and Ottoman navies were already fighting desperately in their home waters; to them the sea was naturally a battlefield and embarking unarmed madness.

This reasoning is not foolish. Its problem is evidence. The Indian Ocean provides a comparison: a thousand masterless years of peaceful prosperity, followed after 1511 by Portuguese control of only a few bases despite enormous effort. Pass fees could not cover military costs; displaced traders shifted to Aceh, Banten, and Johor while Malacca declined. Armed monopoly did not create trade; it taxed existing trade. Our first counterexample exposes the mistake of confusing order's beneficiary with its creator.

\section[Thinking Bridge: Points and Lines]{Thinking Bridge: Change the Map from Blocks to Points and Lines}
Our portable tool changes historical perspective. Territorial maps ask who owns a place; network maps ask what moves through which points. Emperors and armies lead the former; merchants, sailors, goods, languages, and gods lead the latter.

Remember three components wherever you use it.

First: nodes, places where people and goods pause and transfer: ports, oases, post stations, today's airports and data centres. Their value lies less in their size than in crossing links. Small, unproductive Malacca became world-class because Southeast Asian spices, Indian cotton, and Chinese porcelain passed its door. Hubs grow rich through passage, not production. Look at Singapore, or which hometown railway station makes a small city lively.

Second: links and their timetables. Routes have direction, season, and cost. Summer heat over South Asia draws moist sea air inland, producing southwest winds; winter's faster land cooling reverses them. This clock has run for millions of years more punctually than emperors. East Africa to India follows summer winds, the return winter winds. Trading from East Africa all the way to China requires years of stages or handing cargo to specialists along the next stretch. Hence a relay rather than a direct race.

Third: no centre is permanent. Territorial capitals fall and states die; network hubs relocate. Portuguese Malacca sends merchants to Aceh; expensive Aden sends ships to Hormuz. This ability to move is powerful immunity: a network that bypasses you resists conquest. The ocean lacked a master not for lack of ambition but because ownership's costs were high and returns low.

Revisit other histories with these parts. Silk Roads use oasis nodes, caravan links, and seasonal or wartime timetables. The internet does too. Even classroom gossip has hubs among well-connected pupils, with rerouting after friendships break. A tool's portability gives it value.

\section{A Two-Thousand-Year-Old Itinerary}
Read one of official Chinese history's earliest detailed Indian Ocean route accounts, in the Book of Han's geographical treatise. Compiled in the first century CE, it describes earlier Western Han events:

\begin{quote}
From the frontier posts of Rinan, Xuwen, and Hepu, sailing about five months reaches Duyuan; another four months reaches Yilumo... From Fugandulu, over two months' sailing reaches Huangzhi... South of Huangzhi is Yichengbu, from which Han interpreter-envoys returned. ---Book of Han, Treatise on Geography.
\end{quote}
Unpack the itinerary. Rinan, Xuwen, and Hepu were southern Han departure points, roughly today's northern Vietnam and Guangdong and Guangxi coasts. Five months led to Duyuan, four more to Yilumo. Ancient names are difficult to identify precisely, but Huangzhi is generally placed on southeastern India's coast. The same passage says envoys took gold and silk to buy pearls, glass, unusual stones, and other goods. Crucially, ``foreign merchants' ships conveyed them onwards in stages'': overseas, Han envoys changed onto local trading vessels.

This tells us three things.

First, the route is astonishingly old. Around the beginning of the Common Era, travel from southern China to southern India already had established timings, five months and four months, like train schedules. Recorded durations suggest sufficiently frequent travel to become common knowledge.

Second, it belonged to merchants from the start, not emperors. Beyond the border, imperial envoys lacked official ships and relied on foreign traders. No royal road crossed the sea; imperial grandeur stopped at Hepu. If envoys needed merchant transport, ordinary cargo and passengers surely relied on it too.

Third, ``conveyed onwards in stages'' directly evidences a relay network: no single ship from Hepu to Huangzhi, but successive local specialists.

Later records multiplied. Southern Song customs official Zhao Rukuo never travelled abroad but wrote the Description of Foreign Lands from daily conversations with merchants, recording products, routes, and trade across dozens of countries. Yuan traveller Wang Dayuan made two merchant voyages and described over two hundred places from Southeast Asia to East Africa in A Brief Account of Island Barbarians, largely firsthand. Ming interpreter Ma Huan accompanied Zheng He's fleets and in The Overall Survey of the Ocean's Shores vividly described Calicut bargaining: buyer and seller negotiated before a broker, sealed the price by clapping hands, and did not retract the agreement.

Practise weighing evidence: official compilation, secondhand notes, and firsthand travel have different reliability. Even witnesses need scrutiny. Fourteenth-century Ibn Battuta claimed to visit China, but historians still debate actual arrival versus borrowed reports. Travel writing is not surveillance footage. People misremember, boast, and tailor descriptions to home audiences. This does not make all sources untrustworthy; each must earn trust separately.

\section[Order without Ocean Police]{Why Was an Unpoliced Ocean Orderly? Three Explanations}
Return to the core claim: for over a millennium, no unified hegemon ruled the Indian Ocean, yet commerce was broadly peaceful and prosperous. Distinguish four categories. Events: trade's existence and prosperity, supported by shards, wrecks, ports, and texts. Causes: our contested explanations. Contemporary understanding: sultans credit governance, merchants reputation, sailors respect for winds. Evaluation: whether this evidences a better world, reserved for the end.

Compare three genuinely different explanations.

\textbf{Explanation one: the monsoon policed everyone.}

Geography supplied a faultless million-year clock, bringing ships into and out of ports together. Simultaneity matters: rulers had only one or two annual chances to host hundreds of ships witnessing their behaviour; reputation travelled fast. Pirates faced the same schedule, concentrating both targets and protection. Winds coordinated everyone without orders. This explains rhythm, but monsoons preceded flourishing trade by millions of years. A clock alone cannot create a market nobody wishes to attend.

\textbf{Explanation two: merchants brought their own order.}

Institutions and trust deserve credit. Shared faith and origins supported credit, partnerships, arbitration, and collective exclusion of defaulters, harsher than official punishment. Islam's commercial roots mattered: Muhammad's merchant-family origins are well known to believers. Trade carried Islam from Arabia through India into Southeast Asia, where merchants and religious teachers, rather than armies, chiefly spread it. Indonesia, now the largest Muslim-population country, was scarcely ever conquered by an external Muslim state. Language travelled too: Bantu-structured Swahili absorbed Arabic vocabulary, even taking its name from the Arabic for coasts. A language becomes trade history. This explains daily order but hides exclusion: independent newcomers and outsiders lacked network protection. Trust could deter unpaid debts, not organised force; Portuguese cannon in 1511 left reputational networks unable to fight back.

\textbf{Explanation three: fragmentation itself protected trade.}

Counterintuitively, no master produced order. Dozens of coastal states lacked power to levy protection money across the whole ocean, so competed through low taxes, fairness, and quick dispute resolution. Merchants voted with their feet: arbitrary Malacca prices sent next year's ships elsewhere. Monopoly emboldens abuse; competition demands service. This explains prosperity without hegemony, but fragmentation also meant coastal warfare and openings for pirates. Nor does it explain why the functioning system stood nearly undefended before sixteenth-century European armed traders. Competition constrained taxing sultans, not outsiders seeking control of its lifelines rather than its markets.

Monsoons supplied rhythm, merchants daily order, fragmentation freedom. Together they form a fuller cable. The vulnerable ending reminds us that intricate internal order may be unprepared for outsiders. That is our next question.

\section[Further Challenge: Braudel's Longue Durée]{Further Challenge: Make the Ocean the Protagonist---Braudel's Longue Durée}
Introduce French historian Fernand Braudel, one of the twentieth century's most important historians, whose great work was The Mediterranean and the Mediterranean World in the Age of Philip II. Its method travels; later historians, including Indian-born K. N. Chaudhuri, applied it to the Indian Ocean.

Braudel's central claim makes history three superimposed speeds rather than one.

Slowest is geographical time: monsoons, currents, coastlines, the Strait of Malacca's width, Cape storms. Almost motionless over tens of thousands of years, these do not fight wars but determine their possible months.

In the middle lies social time: institutions, routes, religious networks, languages, commercial rules, changing across decades or centuries. Gujarati credit, Swahili's formation, and Southeast Asian Islamisation belong here.

Fastest is event time: battles, treaties, royal deaths, voyages. Traditional histories favour this news-like layer of vivid heroes and villains. Da Gama's 1498 arrival and Malacca's 1511 capture belong here.

Braudel shows that the noisiest layer is shallowest. Great surface waves scarcely move deep water. A sultan can wage war, not redirect monsoons; an emperor can ban shipping, but enduring spice profits and winds send traders out elsewhere. Bans alter prices, not structure. Ming China offers a case: Zheng He's seven early-fifteenth-century voyages projected enormous state power to East Africa, their treasure ships remembered for decades. Then official expeditions abruptly stopped. If emperors wrote all history, Chinese maritime presence should have ended. It did not. Fujian merchants still sailed south to Malacca, and ceramics still decorated East African walls. The emperor left; the network remained. Structures outlive events.

Braudel also proposed the économie-monde, or world-economy: a region forming an integrated trading whole with division of labour, recognised central cities, accepted currency and rules, and a shared rhythm. The Mediterranean was one; the Indian Ocean another, striking for mobile, multiple centres acknowledging no single superior while operating for a millennium. Earlier sections have described this concept in language for younger readers.

Weigh its limits too. The long view resembles a high-flying camera, showing winds and routes but losing faces. Structures may seem to determine everything and reduce individuals to puppets, yet every real decision belongs to someone: a Swahili sailor turning before a storm, a Gujarati merchant trusting three years' profit to a distant partner. No structure prewrites every choice. It provides a stage and draft script, not performers. Beware theories turning everything into inevitability. One of history's most valuable uncertainties is that alternatives existed then.

\section{Today's Monsoon Still Blows}
The old problem remains, differently dressed.

The Strait of Malacca is still among the world's busiest waterways, now carrying tankers and vessels with tens of thousands of containers. Monsoons no longer schedule departures, but precise shipping rhythms persist. Parts of your phone, trainers, and games console probably crossed these waters. The spice-price logic now appears as global division of labour: design in one country, components in three, assembly in another, much of the price reflecting passage. Singapore, a small island producing almost nothing, grew into a great port and financial centre at the strait's exit: Malacca's direct descendant.

The old question returns: who governs the sea? Great-power navies patrol and display control; Hormuz tanker passages unsettle oil prices; freedom-of-navigation disputes recur. Separate facts, such as verifiable ship counts, from explanations contesting order versus monopoly, and from your judgement about maritime policing. A millennium's counterintuitive answer, order need not require ownership, may not fit today. It still challenges land-based assumptions that masterlessness automatically means chaos.

Quieter echoes survive. Malaysia's Baba-Nyonya communities descend from Chinese merchants who settled with the monsoons and married local women centuries ago. Malay language, local clothing, and cuisine combining Chinese methods and local spices make them living descendants of our port families. Swahili is now spoken daily by tens of millions, from Kenyan classrooms to Tanzania's parliament. Languages and communities outlive empires: sultans are gone, mixed cultures endure.

\section{Your Question: Does the Sea Need a Master?}
Return to the title: no master, never empty.

Which is better, a sea with a master or without? Answer with reasons.

First weigh both sides again.

For a master: order is a public good someone must finance. Without police, weak ships suffer first; rich merchants hire guards while small traders accept their fate. You have reconstructed Portugal's serious case, and much order you enjoy today really is maintained at someone's expense.

Against a master: monopoly enables abuse, competition requires service. The masterless millennium sustained more languages and faiths than any empire, leaving Baba-Nyonya and Swahili as living monuments. Would-be master Portugal left only forts and military deficits.

A harder layer remains. The masterless golden age welcomed insiders but shut out strangers. Was its beauty only a members' club's beauty? Meanwhile, the master's harsh order issued tickets to everyone.

What do you choose? Or is there a missing third possibility: a sea monopolised by no empire but supported jointly by all?

There is no standard answer. Next time an atlas shows a broad blank blue ocean, may you see wind, ports, and people waiting for wind, not emptiness. It was never empty.
